\subsubsection{GLUE accuracy under FlashSVD v1.5}

Table~\ref{tab:glue_main_flashsvd15} reports GLUE results under the \texttt{flashsvd15} backend for the dense baseline, compressed models, and compressed models after post-compression fine-tuning.
Overall, post-compression fine-tuning substantially recovers task accuracy across all methods, with DRONE achieving the strongest A-GLUE and A-Acc among the compressed+FT variants.
Figure~\ref{fig:glue_combined_main} provides a compact visual summary of the same trend, showing that stage-1 compression causes noticeable accuracy degradation while stage-2 post-compression fine-tuning largely restores performance across all methods.

\begin{table*}[t]
\centering
\small
\renewcommand{\arraystretch}{1.10}
\setlength{\tabcolsep}{2pt}
\begin{threeparttable}
\caption{
GLUE results under the \texttt{flashsvd15} backend for the dense baseline, compressed models, and compressed models after post-compression fine-tuning.
All numbers are percentages; higher is better.
}
\label{tab:glue_main_flashsvd15}
\begin{tabular}{llcccccccccc}
\toprule
Group & Method & CoLA & SST-2 & MRPC & QQP & MNLI & QNLI & RTE & STS-B & A-GLUE & A-Acc \\
\midrule
Original
& BERT$_{\mathrm{base}}$
& 53.4 & 92.4 & 91.0 & 87.8 & 84.6 & 91.6 & 72.6 & 88.1 & 82.7 & 85.3 \\
\midrule
\multirow{4}{*}{Compressed}
& SVD
& 1.7 & 71.2 & 0.0 & 16.7 & 37.4 & 54.7 & 47.3 & 35.3 & 33.0 & 52.6 \\
& FWSVD
& 18.9 & 72.0 & 10.1 & 67.9 & 52.4 & 53.2 & 58.1 & 70.2 & 50.4 & 58.9 \\
& DRONE
& 1.9 & 85.3 & 84.8 & 76.2 & 57.9 & 60.2 & 59.9 & 49.4 & \textbf{59.5} & \textbf{65.8} \\
& AdaSVD
& -1.5 & 62.0 & 0.7 & 52.1 & 35.2 & 45.9 & 54.9 & 63.3 & 39.1 & 49.5 \\
\midrule
\multirow{4}{*}{Compressed + FT}
& SVD
& 28.3 & \textbf{91.1} & 84.8 & 87.4 & 82.7 & 88.8 & 59.6 & 86.4 & 76.1 & 80.5 \\
& FWSVD
& \textbf{47.5} & 90.5 & 86.7 & \textbf{87.8} & 82.8 & 89.3 & 65.7 & 87.0 & \textbf{79.7} & 82.1 \\
& DRONE
& 42.1 & 90.8 & \textbf{89.7} & 87.4 & 82.4 & 89.4 & \textbf{70.4} & 85.4 & \textbf{79.7} & \textbf{83.3} \\
& AdaSVD
& 39.4 & \textbf{91.1} & 89.5 & 87.6 & \textbf{82.5} & \textbf{89.8} & 62.8 & \textbf{87.1} & 78.7 & 81.5 \\
\bottomrule
\end{tabular}
\begin{tablenotes}[flushleft]
\footnotesize
\item CoLA is reported with Matthews correlation, MRPC and QQP with F1, STS-B with Pearson correlation, and SST-2/MNLI/QNLI/RTE with accuracy.
\item A-GLUE averages all eight task metrics as reported in the table.
\item A-Acc averages only the accuracy-based tasks: SST-2, MNLI, QNLI, and RTE.
\item Boldface highlights the best compressed result within each compression stage.
\end{tablenotes}
\end{threeparttable}
\end{table*}


\begin{figure}[t]
    \centering
    \includegraphics[width=0.95\linewidth]{figures/encoders/accuracy/glue_both_stages.png}
    \caption{
    GLUE average comparison across compression methods before and after post-compression fine-tuning.
    Results are shown at bf16 for the per-head compression setting used in the main encoder experiments.
    The left panel reports stage-1 performance immediately after compression, while the right panel reports stage-2 performance after post-compression fine-tuning.
    The figure shows that stage-1 accuracy drops vary substantially across methods, whereas fine-tuning recovers most of the lost performance and narrows the gap to the dense baseline.
    Higher is better.
    }
    \label{fig:glue_combined_main}
\end{figure}

\subsubsection{Backend efficiency on real GLUE inputs}

Table~\ref{tab:backend_comparison_glue} summarizes backend efficiency averaged over compressed GLUE runs.
Among the evaluated backends, \texttt{flashsvd15} delivers the best latency and throughput, while \texttt{flashsvd} attains slightly lower peak memory.
Figures~\ref{fig:throughput_backend_main} and~\ref{fig:memory_backend_main} show the same trend from the throughput and memory perspectives, respectively: FlashSVD v1.5 consistently achieves the strongest throughput, while both FlashSVD runtimes remain far more memory-efficient than the naive and SDPA baselines.

\begin{table}[t]
\centering
\small
\setlength{\tabcolsep}{4pt}
\renewcommand{\arraystretch}{1.10}
\begin{threeparttable}
\caption{
Backend comparison on GLUE inference efficiency for compressed BERT models.
Results are averaged over SVD, FWSVD, DRONE, and AdaSVD, and over all eight GLUE tasks using real inputs.
Lower latency and memory are better; higher throughput and speedup are better.
}
\label{tab:backend_comparison_glue}
\begin{tabular}{lcccc}
\toprule
Backend & Avg. Latency (ms) $\downarrow$ & Avg. Throughput $\uparrow$ & Peak Mem. (MB) $\downarrow$ & Speedup vs naive $\uparrow$ \\
\midrule
naive      & 64.99 & 492.4  & 1007.2 & 1.00$\times$ \\
sdpa       & 26.61 & 1203.1 & 623.2  & 2.44$\times$ \\
flashsvd   & 44.20 & 724.0  & \textbf{359.8} & 1.47$\times$ \\
flashsvd15 & \textbf{21.79} & \textbf{1471.1} & 361.7 & \textbf{2.98$\times$} \\
\bottomrule
\end{tabular}
\begin{tablenotes}[flushleft]
\footnotesize
\item Results are averaged over CoLA, SST-2, MRPC, QQP, MNLI, QNLI, RTE, and STS-B.
\item Results are further averaged over SVD, FWSVD, DRONE, and AdaSVD.
\item Speedup vs naive is computed from average latency.
\end{tablenotes}
\end{threeparttable}
\end{table}


\begin{figure}[t]
    \centering
    \includegraphics[width=0.92\linewidth]{figures/encoders/backend_compare/throughput_mnli.png}
    \caption{
    Throughput comparison on representative compressed checkpoints under different backends.
    Results are shown at bf16, sequence length 512, and batch size 32 for representative GLUE tasks.
    We compare \texttt{sdpa}, \texttt{flashsvd}, and \texttt{flashsvd15} across compressed methods.
    FlashSVD v1.5 consistently achieves the highest throughput, substantially outperforming FlashSVD v1 and the naive baseline, and often exceeding SDPA as well.
    Higher is better.
    }
    \label{fig:throughput_backend_main}
\end{figure}

\begin{figure}[t]
    \centering
    \includegraphics[width=0.92\linewidth]{figures/encoders/backend_compare/peak_memory_mnli.png}
    \caption{
    Peak GPU memory comparison on representative compressed checkpoints under different backends.
    Results are shown at bf16, sequence length 512, and batch size 32 for representative GLUE tasks.
    The figure compares backend-dependent memory usage across compressed methods, showing that \texttt{flashsvd} attains slightly lower peak memory while \texttt{flashsvd15} remains close and far below the naive and SDPA baselines.
    Lower is better.
    }
    \label{fig:memory_backend_main}
\end{figure}

\subsubsection{Memory breakdown and scaling behavior}

To better understand where the memory gains come from, we decompose peak inference memory into parameter storage and activation-related overhead.
Figure~\ref{fig:memory_breakdown_main} shows that low-rank compression reduces parameter memory, but the naive low-rank backend still suffers from large activation and buffer overhead.
In contrast, FlashSVD eliminates most of this overhead, which is why the total peak memory can fall even below the dense baseline.

We further study whether these gains persist under larger workloads.
Figure~\ref{fig:seqlen_scaling_main} shows sequence-length scaling, and Figure~\ref{fig:batch_scaling_main} shows batch-size scaling.
Across both dimensions, FlashSVD v1.5 maintains the best throughput while preserving the memory advantage of fused low-rank execution, indicating that its benefits are not limited to a single operating point.

\begin{figure}[t]
    \centering
    \includegraphics[width=0.92\linewidth]{figures/encoders/memory_analysis/memory_breakdown.png}
    \caption{
    Memory breakdown into parameter storage and activation-related overhead for dense and compressed encoder inference.
    Results are shown at bf16, sequence length 512, and batch size 32.
    The figure illustrates that compression reduces parameter memory, while the FlashSVD backend further reduces total peak memory by eliminating activation and buffer overhead that remains large under naive low-rank execution.
    }
    \label{fig:memory_breakdown_main}
\end{figure}

\begin{figure*}[t]
    \centering
    \includegraphics[width=0.95\linewidth]{figures/encoders/seqlen_scaling/all_panels.png}
    \caption{
    Sequence-length scaling for SVD-compressed BERT under different backends.
    From left to right, we report peak GPU memory, throughput, and relative memory reduction versus the naive einsum backend.
    Results are collected at bf16 with batch size 32 using synthetic full-length inputs, and are averaged over compressed methods and tasks.
    The figure compares how backend efficiency changes as sequence length increases from 128 to 512.
    FlashSVD backends scale more favorably with sequence length than naive and SDPA baselines, with FlashSVD v1.5 achieving the highest throughput and the largest relative memory reduction at longer sequence lengths.
    }
    \label{fig:seqlen_scaling_main}
\end{figure*}

\begin{figure*}[t]
    \centering
    \includegraphics[width=0.95\linewidth]{figures/encoders/batch_scaling/all_panels.png}
    \caption{
    Batch-size scaling for SVD-compressed BERT under different backends.
    The left panel reports peak GPU memory, and the right panel reports throughput.
    Results are collected at bf16 with sequence length 512 using synthetic full-length inputs, and are averaged over compressed methods and tasks.
    We compare the naive einsum backend, SDPA, FlashSVD v1, and FlashSVD v1.5 as batch size increases from 8 to 64.
    The figure compares how backend choice affects memory growth and throughput scaling under increasing workload parallelism.
    FlashSVD backends substantially reduce memory across all batch sizes, while FlashSVD v1.5 delivers the strongest throughput, especially at moderate and large batch sizes.
    }
    \label{fig:batch_scaling_main}
\end{figure*}

\subsubsection{Real vs. synthetic benchmarking}

We additionally compare backend speedups measured on real GLUE inputs and on synthetic full-length inputs.
As shown in Table~\ref{tab:real_vs_synth_gap}, speedups on real GLUE inputs are consistently higher than those measured under synthetic full-length inputs.
This pattern is especially pronounced for \texttt{flashsvd15}, indicating that synthetic full-length benchmarking can understate realized end-to-end gains on the real GLUE workloads considered here.

\begin{table}[t]
\centering
\small
\setlength{\tabcolsep}{29pt}
\renewcommand{\arraystretch}{1.08}
\begin{threeparttable}
\caption{
Real vs. synthetic speedup on GLUE workloads at sequence length 512 and batch size 32.
We report backend speedup relative to the naive backend.
}
\label{tab:real_vs_synth_gap}
\begin{tabular}{lccc}
\toprule
Backend & Real $\uparrow$ & Synthetic $\uparrow$ & Difference \\
\midrule
SDPA          & 2.44$\times$ & 2.01$\times$ & 0.43 \\
FlashSVD v1   & 1.47$\times$ & 1.15$\times$ & 0.32 \\
FlashSVD v1.5 & 3.02$\times$ & 2.35$\times$ & 0.67 \\
\bottomrule
\end{tabular}
\begin{tablenotes}[flushleft]
\footnotesize
\item Values are averaged over compressed GLUE runs (SVD, FWSVD, DRONE, and AdaSVD; eight tasks).
\item Real GLUE inputs contain substantial sequence padding, while synthetic inputs use fixed full-length 512-token sequences.
\item Difference is computed as Real minus Synthetic.
\item Larger differences indicate a stronger gap between realized speedup on real workloads and speedup measured under synthetic full-length inputs.
\end{tablenotes}
\end{threeparttable}
\end{table}


\subsubsection{Kernel-level analysis}

To better understand the efficiency differences, we inspect kernel-level execution on MNLI.
Table~\ref{tab:kernel_profile_mnli} shows that fused FlashSVD backends reduce execution diversity and simplify the kernel mix relative to the naive path.
This effect is especially clear for AdaSVD under \texttt{flashsvd15}, which exhibits fewer kernel calls, fewer GEMM variants, and lower fragmentation.
Figure~\ref{fig:kernel_breakdown_main} complements this view by showing that heterogeneous-rank compression increases kernel diversity in the naive path, while FlashSVD backends consolidate dispatch and shift a substantial fraction of runtime into fused execution.

\begin{table}[t]
\centering
\scriptsize
\setlength{\tabcolsep}{9pt}
\renewcommand{\arraystretch}{1.10}
\begin{threeparttable}
\caption{
Kernel-profile summary on MNLI for representative compressed BERT checkpoints.
Results are collected at bf16 and sequence length 512.
We compare the naive, \texttt{flashsvd}, and \texttt{flashsvd15} backends in terms of total kernel calls, GEMM kernel diversity, fragmentation, and the dominant kernel by time share.
Lower kernel counts, fewer GEMM variants, and lower fragmentation indicate a more streamlined execution pattern.
}
\label{tab:kernel_profile_mnli}
\begin{tabular}{llrrrr}
\toprule
Method & Backend & Kernel Calls $\downarrow$ & GEMM Variants $\downarrow$ & Fragmentation $\downarrow$ & Top Kernel (Time \%) \\
\midrule
SVD
& naive      & 23000 & 8  & 920.0 & GEMM (17.1) \\
& flashsvd   & \textbf{12307} & \textbf{5} & \textbf{559.4} & fused\_ffn\_phase1 (49.4) \\
& flashsvd15 & 21200 & \textbf{5} & 963.6 & fused\_ffn\_full\_batched\_bias (37.5) \\
\midrule
AdaSVD
& naive      & 23000 & 14 & 741.9 & softmax\_warp\_forward (13.3) \\
& flashsvd   & 19400 & 10 & 718.5 & fused\_ffn\_phase1 (50.2) \\
& flashsvd15 & \textbf{14000} & \textbf{9} & \textbf{560.0} & fused\_ffn\_full\_batched\_bias (42.4) \\
\bottomrule
\end{tabular}
\begin{tablenotes}[flushleft]
\footnotesize
\item Results are collected from MNLI profiling runs.
\item FlashSVD v1 already reduces execution diversity for SVD checkpoints.
\item FlashSVD v1.5 shows the clearest kernel-streamlining effect on AdaSVD.
\end{tablenotes}
\end{threeparttable}
\end{table}


\begin{figure*}[t]
    \centering
    \includegraphics[width=0.95\linewidth]{figures/encoders/nsys/combined.png}
    \caption{
    Kernel-level utilization analysis under uniform-rank and heterogeneous-rank compression on MNLI.
    Results are collected at bf16 and sequence length 512.
    The left panel reports the number of distinct GEMM kernel configurations, and the right panel shows GPU time breakdown across fused Triton kernels, cuBLAS/CUTLASS GEMM kernels, and other operations.
    We compare SVD and AdaSVD checkpoints under the naive backend, FlashSVD v1, and FlashSVD v1.5.
    The figure is designed to compare how rank heterogeneity changes execution structure and how much of the runtime can be shifted into fused kernels.
    Heterogeneous-rank compression increases kernel diversity under the naive backend, while both FlashSVD backends consolidate dispatch and shift a substantial fraction of runtime into fused execution.
    }
    \label{fig:kernel_breakdown_main}
\end{figure*}
